% 15 — De Morgan's laws for sets (two rows of identical shaded mini-Venns)
\begin{tikzpicture}[font=\itshape]
  % mini-Venn: complement of union (shade outside both circles)
  \newcommand{\vUC}[3]{\begin{scope}[shift={(#1,#2)}]
      \coordinate (a) at (-0.45,0);\coordinate (b) at (0.45,0);\def\r{0.85}
      \begin{scope}\clip (-1.6,-1.2) rectangle (1.6,1.2);
        \fill[#3,opacity=0.4] (-1.6,-1.2) rectangle (1.6,1.2);
        \fill[white] (a) circle(\r) (b) circle(\r);\end{scope}
      \draw[gray,line width=.8pt] (-1.6,-1.2) rectangle (1.6,1.2);
      \draw[navy,line width=.8pt] (a) circle(\r);
      \draw[Fred,line width=.8pt] (b) circle(\r);
      \node[navy,font=\scriptsize\itshape] at (-1.05,.95){$A$};
      \node[Fred,font=\scriptsize\itshape] at (1.05,.95){$B$};\end{scope}}
  % mini-Venn: complement of intersection (shade all except lens)
  \newcommand{\vIC}[3]{\begin{scope}[shift={(#1,#2)}]
      \coordinate (a) at (-0.45,0);\coordinate (b) at (0.45,0);\def\r{0.85}
      \begin{scope}\clip (-1.6,-1.2) rectangle (1.6,1.2);
        \fill[#3,opacity=0.4] (-1.6,-1.2) rectangle (1.6,1.2);
        \begin{scope}\clip (a) circle(\r);\fill[white](b) circle(\r);\end{scope}\end{scope}
      \draw[gray,line width=.8pt] (-1.6,-1.2) rectangle (1.6,1.2);
      \draw[navy,line width=.8pt] (a) circle(\r);
      \draw[Fred,line width=.8pt] (b) circle(\r);
      \node[navy,font=\scriptsize\itshape] at (-1.05,.95){$A$};
      \node[Fred,font=\scriptsize\itshape] at (1.05,.95){$B$};\end{scope}}
  % ROW 1
  \node[formula,font=\itshape\large] at (0,3.2) {$(A \cup B)^{c} = A^{c} \cap B^{c}$};
  \vUC{-2.6}{1.5}{orange}
  \node[gray,font=\Large\upshape] at (0,1.5) {$=$};
  \vUC{2.6}{1.5}{orange}
  % ROW 2
  \node[formula,font=\itshape\large] at (0,-0.4) {$(A \cap B)^{c} = A^{c} \cup B^{c}$};
  \vIC{-2.6}{-2.1}{sgreen}
  \node[gray,font=\Large\upshape] at (0,-2.1) {$=$};
  \vIC{2.6}{-2.1}{sgreen}
\end{tikzpicture}
